\section{Q\&A 延伸：课堂 slide 没写出的限制}

Q\&A 没有新增独立 visual state，却补上了几条决定技术结论是否可靠的限定：long context 的 prefill cost、data 与 architecture 的非零和关系、diffusion 优势的不确定性、CogAgent 的高分辨率来源，以及视频能否从 raw signal 学到物理。

\subsection{Long context、data 与 architecture 的三条平衡线}

第一，长上下文可减少外部 retrieval pipeline，但 prefill latency 与无关 token 干扰仍然存在。第二，data 可以表达许多 task-specific biases，但 general architecture improvement 仍可能改变所有任务的 compute frontier。第三，把全部 documents 或 reasoning branches 放入 context 需要模型学会利用，而不是只扩大窗口。

\begin{importantbox}{课堂 Q\&A 的工程结论}
\begin{enumerate}
  \item 对长文档短回答，允许较慢 prefill 可能合理；实时 agent 则未必。
  \item 对具体行为，先尝试 data/evaluation；对跨任务重复瓶颈，再考虑 architecture。
  \item 若 reasoning procedure 包含失败路径，把它们显式放入 context 可能比只保留 beam score 更有信息，但仍需实验。
\end{enumerate}
\end{importantbox}

\teachervoice{01:18:10--01:19:59，讲者倾向把 alternative paths 与 failure information 一并放入 context，让模型学习比较，而不是只用 hard-coded beam probability；这是一条经验，不是 universal theorem。}

\subsection{Video understanding 不会自动获得物理世界}

当前 VLM/video model 多依赖 text-image 或 text-video pairs。若 training signal 主要来自人工 annotation，模型学到的是人类描述覆盖的现象，不必然从 raw motion 中恢复 conservation、contact、object permanence 或 causal interventions。

\begin{warningbox}{“看过很多视频”不等于“学会物理”}
要从未标注视频学习物理，可能需要 temporal prediction、masked video modeling、world-model objectives、interaction data 或可验证 simulation。仅扩大 annotated pairs，仍可能学到语言 shortcut 与表面相关性。
\end{warningbox}

\teachervoice{01:16:15--01:18:10，讲者明确说如果数据源不包含 physical rules，现有 pair-based pretraining 不会自动学出它们；他把新的 video self-supervision 视为必要研究问题。}

\subsection{本章小结}

Q\&A 把主讲中的 headline 还原成带条件的工程判断：context 有延迟，data 不能消灭 architecture，diffusion 没有数学上击败 autoregression，GUI demo 不代表可靠 agent，video quantity 也不等于 physical understanding。

